\begin{hongqiarticle}{
        title  = {坚持以阶级斗争为纲},
        author = {程越},
        style  = normal,
        kicker = {努力学习毛泽东思想 批判邓小平的修正主义路线}
    }

    在全党全军全国各族人民怀着极其悲痛的心情哀悼伟大的领袖和导师毛泽东
    主席逝世的时候，我们不会忘记：毛主席在批邓、反击右倾翻案风的斗争中，又
    一次谆谆教导广大党员、干部要抓好阶级斗争，并指出“阶级斗争是纲，其余都
    是目”。我们一定要继承毛主席的遗志，坚持以阶级斗争为纲，坚持党的基本路
    线，坚持无产阶级专政下的继续革命，深入批邓、反击右倾翻案风，把无产阶级
    革命事业进行到底。

    几千年来的阶级社会的历史表明，阶级斗争是推动社会向前发展的强大动力。
    坚持以阶级斗争为纲，符合历史的发展规律，符合马克思列宁主义的基本原理。抓
    住无产阶级同资产阶级的阶级斗争，社会主义同资本主义两条道路斗争这个纲，是
    毛主席领导全国人民进行社会主义革命和建设的基本理论和基本实践。中华人民
    共和国成立前夕，毛主席就明确指出，社会主义革命时期国内的主要矛盾是“工人
    阶级和资产阶级的矛盾”。（《在中国共产党第七届中央委员会第二次全体会议上的报告》）
    毛主席这一重要指示，为我们同党内外资产阶级作斗争奠定了坚实的理论基础。在
    生产资料所有制的社会主义改造基本完成后，事实上还存在着阶级和阶级斗争，而
    刘少奇一伙党内资产阶级为了向无产阶级进攻，却极力否认这一点。毛主席极其
    深刻地批驳了刘少奇、邓小平等散布的否认阶级和阶级斗争的谬论，指出：“被推翻
    的地主买办阶级的残余还是存在，资产阶级还是存在，小资产阶级刚刚在改造。
    阶级斗争井没有结束。无产阶级和资产阶级之间的阶级斗争，各派政治力量之间
    的阶级斗争，无产阶级和资产阶级之间在意识形态方面的阶级斗争，还是长时期
    的，曲折的，有时甚至是很激烈的”。（《关于正确处理人民内部矛盾的问题》）经过十三
    年的社会主义革命和建设的伟大实践，针对国内外阶级敌人的疯狂进攻，针对党
    内修正主义路线头子的阴谋活动，毛主席在一九六二年发出了“千万不要忘记阶
    级和阶级斗争”的伟大号召，并且更加完整地提出了党在整个社会主义历史阶段
    的基本路线，教导我们对社会主义时期存在着阶级、阶级矛盾和阶级斗争的问题，
    “必须年年讲，月月讲，天天讲”，使大家有比较清醒的认识。随后，在城乡开展
    的社会主义教育运动中，在无产阶级文化大革命中，毛主席对社会主义时期阶级
    斗争的新形势，新特点，又作出了科学的深刻的分析，明确指出走资本主义道路
    的当权派是社会主义革命的重点对象。在批邓、反击右倾翻案风的斗争中，毛主席
    进一步作出了“搞社会主义革命，不知道资产阶级在哪里，就在共产党内，党内走
    资本主义道路的当权派。走资派还在走”的科学论断。在深切悼念毛主席逝世的
    时候，我们重温毛主席关于以阶级斗争为纲的光辉教导，感到多么亲切、多么重
    要啊！

    毛主席关于以阶级斗争为纲的一系列指示，是对国际和国内马克思主义同修
    正主义、无产阶级同资产阶级之间长期激烈的斗争所作出的科学总结，大大丰
    富了马克思主义的理论宝库，是我们坚持无产阶级专政下的继续革命，反修防
    修，巩固无产阶级专政，防止资本主义复辟，建设社会主义的无价之宝。建国二
    十多年来，从“三反”“五反”、社会主义改造，反右派斗争到社会主义教育运动，
    从揭露高岗饶漱石反党联盟，彭德怀反党集团到文化大革命中摧毁刘少奇、林彪
    两个资产阶级司令部，从批判邓小平反革命的修正主义路线到粉碎天安门广场的
    反革命政治事件，都是毛主席领导我们同党内外资产阶级所作的一系列极其尖锐、
    极其激烈的斗争。随着每一次斗争的胜利，社会主义事业都大大向前推进了一步。
    我们的党，我们无产阶级专政的社会主义国家，正是在这种阶级斗争中得到巩固
    和加强的。

    在社会主义历史阶段要不要坚持以阶级斗争为纲，一直是毛主席的无产阶级
    革命路线同反革命的修正主义路线进行激烈斗争的焦点，是区别真假马克思主义
    的重要标志。在我们党内出现的修正主义路线的头子高岗，饶漱石、彭德怀、刘
    少奇，林彪、邓小平，同老的社会民主党和以苏修叛徒集团为中心的现代修正主
    义一样，都把阶级斗争熄灭论作为他们的路线和政策的理论基础。他们根本否认
    几千年来的历史是阶级斗争的历史，根本否认无产阶级对资产阶级的斗争，根本
    否认无产阶级对资产阶级的专政，猖狂反对我们党在社会主义时期的基本理论和
    基本实践。每当社会主义革命深入一步，这些修正主义路线的头子总是不择手段
    地反对和歪曲毛主席关于以阶级斗争为纲的指示，大肆散布阶级斗争熄灭论，用
    来麻痹人民群众的革命意志，以掩盖他们向党向社会主义进攻。事实上，阶级斗
    争从未止息，刘少奇大搞资本主义复辟，林彪发动反革命武装政变，邓小平刮起
    右倾翻案风，这些修正主义路线的头子反对革命的行径，就完全截穿了阶级斗争
    熄灭论的骗局。在毛主席关于以阶级斗争为纲的指示深入人心，反动的阶级斗争
    熄灭论被批判得体无完肤的时候，邓小平竞然抛出“三项指示为纲”的反动政治
    纲领，还疯狂地叫嚷：“阶级斗争哪能天天讲？”这正说明邓小平的的确确是不肯改
    悔的走资派，反映了他对无产阶级向资产阶级的斗争既极端仇视、又极端恐惧的
    反动阶级本能。我们天天讲阶级斗争，天天搞阶级斗争，邓小平及其代表的党内
    外资产阶级的日子就不好过，他们的修正主义路线就会受到人民群众的揭露和批
    判，他们复辟资本主义的阴谋活动就必然遭到可耻的失败。邓小平的反动叫嚣从
    反面教育了我们：一定要紧紧抓住阶级斗争这个纲，必须天天抓，月月抓，年年
    抓，任何时候也不可放松。

    坚持以阶级斗争为纲，最重要的是同党内资产阶级作斗争，特别是同党内走
    资派作斗争。刘少奇、林彪、邓小平这些坚持走资本主义道路的当权派，反映了被
    打倒的地主资产阶级和新产生的资产阶级分子的利益和愿望，成为他们的政治代
    表。这些走资派由于窃据了党和国家的部分权力，搞起资本主义复辟来很容易。因
    此，走资本主义道路的当权派对于无产阶级专政的社会主义制度具有极大的危害
    性，是社会主义革命的主要对象、重点对象。如果我们抓不住这个重点，抓阶级斗
    争就会落空。走资派既有很大的反动性，又有很大的欺骗性。要同他们作斗争，就
    要善于识别他们的反革命面目。毛主席提出的“要搞马克思主义，不要搞修正主
    义；要团结，不要分裂；要光明正大，不要搞阴谋诡计”的三项基本原则，是我们识
    别党内走资派、分清真假马克思主义的锐利武器。走资派还在走，这是不以人们的
    意志为转移的。搞修正主义，搞分裂，搞阴谋诡计，这是一切走资派反对无产阶级革
    命事业最本质的特征。只要我们努力学会运用马克思主义、列宁主义、毛泽东思想
    这个政治上的望远镜和显微镜，不管走资派怎样进行伪装，终究会现出他们的原
    形。共产党是领导无产阶级和革命群众对于阶级敌人进行战斗的先锋队。我们共
    产党人的哲学是斗争哲学。作为一个在无产阶级专政下继续革命的中国共产党党
    员，就应该把同资产阶级、特别是同党内资产阶级作斗争作为自已的毕生任务。
    我们要遵循毛主席的教导，响应党中央的号召，誓为反对党内外资产阶级和一切
    剥削阶级，为在地球上消灭一切人剥削人的制度，为实现共产主义而奋斗终生。

    毛主席瞩附我们：“按既定方针办”。正在深入开展的批邓、反击右倾翻案风
    的斗争，是毛主席亲自发动和领导的。这场斗争是反对党内资产阶级的一场伟大
    斗争。我们要把毛主席开创的无产阶级革命事业进行到底，在当前就一定要按
    毛主席既定的部署和政策，进一步搞好批邓。邓小平继承了刘少奇、林彪的衣钵，
    妄图颠覆无产阶级专政，复辟资本主义。邓小平的修正主义路线的反动内容和实
    质，集中反映在他授意炮制的《论总纲》、《汇报提纲》、《条例》这三株大毒草中。
    我们集中火力批邓，就要抓住他搞修正主义的思想政治路线这个要害，继续深
    入批判。邓小平反革命的修正主义路线曾经流毒到各个领域，有它的阶级基础
    和一定的市场。我们不能满足于已经取得的成绩，要乘胜前进。有一些单位重新
    学习了毛主席在反击右倾翻案风斗争中的重要指示，从总结回顾，调查研究入
    手，列专题，抓要点，并分析当前阶级斗争的状况，作出继续深入批邓的规划；
    有一些单位的理论队伍把毛主席继承、捍卫、发展马克思列宁主义的丰功伟绩，同
    邓小平反对马克思主义、列宁主义、毛泽东思想的修正主义纲领和谬论加以比较，
    向群众作了宣讲，使人们更深刻地认识修正主义的反动性和危害性：有一些单位把
    学习中共中央、人大常委会、国务院、中央军委《告全党全军全国各族人民书》和
    华国锋同志在伟大的领袖和导师毛泽东主席追掉大会上致的悼词同批邓结合起来，
    从思想政治路线上把我们如何从各方面继承毛主席的遗志同邓小平如何从各方面
    反对毛主席的革命路线加以对照，有力地推动了批邓的深入。我们要注意总结这些
    来自群众的经验。马克思主义、列宁主义、毛泽东思想是最锐利的批判武器，只有认
    真学好马列著作和毛主席著作，才能把批邓引向深入。各级党组织要加强对学习的
    领导，对于如何进一步学好毛主席著作，要有一个学习计划。学习、批邓都要联系各
    条战线阶级斗争和两条路线斗争的实际，这样才能充分揭露邓小平的修正主义路
    线的反动本质，肃清其流毒和影响，才能更好地提高我们的阶级斗争和路线斗争的
    觉悟，坚决地贯彻执行毛主席的革命路线。对于党内资产阶级和社会上阶级敌人的
    破坏和捣乱，我们一定要提高警惕。对于一小撮阶级敌人散布谣言，制造混乱，妄图
    为邓小平翻案的反革命活动，要充分依靠和发动人民群众予以彻底揭露，坚决打击。

    当前摆在我们面前的任务是十分艰巨的。在繁重的工作中，我们一定要牢记
    毛主席的教导，抓紧阶级斗争这个纲，要用批邓、反击右倾翻案风的斗争，来推
    动各项工作，推动生产。抓住阶级斗争这个纲，就能抵制和批判资本主义、修正
    主义，坚持社会主义道路，就能继续调动亿万人民群众的社会主义积极性，就能
    遵循事物发展的客观规律，高屋建瓴、势如破竹地解决各种问题，搞好各项工
    作。毛主席教导我们：“阶级斗争，一抓就灵。”在悼念伟大的领袖和导师毛主席
    逝世的极其悲痛的时刻，我们要更加牢记毛主席这个教导，紧紧抓住阶级斗争这
    个纲，真正做到抓革命，促生产，促工作，促战备，把各项工作做得更好。

    我们敬爱的伟大领袖和导师毛主席和我们永别了。但是，伟大的毛泽东思想
    将永远照耀我们前进的道路，毛主席的无产阶级革命路线将永远指引我们战斗的
    航程。在党中央的领导下，紧紧“抓住阶级斗争这个纲，抓住社会主义和资本主义两
    条道路斗争这个纲”，我们就一定能够夺取社会主义革命和建设事业的更大胜利！
\end{hongqiarticle}